\begin{table}[htbp]
\centering
\renewcommand{\arraystretch}{1.35}
\small
\caption{تقسیم گروه‌محور استاندارد جمعیت موارد کامل}
\label{tab:split}
\setlength{\tabcolsep}{3pt}
\begin{tabular}{|l|r|r|r|r|r|}
\hline
تقسیم & تعداد & \lr{Normal} & \lr{Pre-Plus} & \lr{Plus} & گروه \\
\hline
آموزش & 6203 & 4490 & 650 & 1063 & 290 \\
\hline
اعتبارسنجی & 1328 & 981 & 140 & 207 & 62 \\
\hline
آزمون & 1331 & 982 & 140 & 209 & 62 \\
\hline
\textbf{مجموع} & \textbf{8862} & \textbf{6453} & \textbf{930} & \textbf{1479} & \textbf{414} \\
\hline
\multicolumn{6}{p{0.96\textwidth}}{\footnotesize هیچ گروهی در بیش از یک تقسیم قرار ندارد. تعریف گروه در منبع \lr{Plus} و فارفوم بر پایه‌ی بیمار و در فارابی بر پایه‌ی معاینه است، زیرا شناسه‌ی بیمار در فارابی در دسترس نبود.} \\
\hline
\end{tabular}
\end{table}